\begin{lemma}[Controlled interchange for a minimum bad pair]
Let \((\varphi,T)\) be a minimum bad pair, with
\(T=(y_0,e_1,y_1,\ldots,e_n,y_n)\).  Suppose \(i<j<n\), the distinct
colors \(\alpha,\beta\) are missing at \(y_i,y_j\), respectively, and
neither color is used on \(e_2,\ldots,e_j\).  Let \(Q\) be the
\((\alpha,\beta)\)-component containing \(y_j\), and obtain
\(\varphi'\) by interchanging the two colors on \(Q\).  Then:
\begin{enumerate}
\item \(Q\) is a path component containing \(y_i\), with \(y_i,y_j\) as
its two vertex ends;
\item \(\beta\) is missing at \(y_i\) and \(\alpha\) at \(y_j\) under
\(\varphi'\);
\item every edge \(e_2,\ldots,e_j\) keeps its exact color, and every old
Tashkinov witness for every later edge gives a (possibly exchanged)
earlier witness for that edge after the interchange; and
\item \((\varphi',T)\) is again a minimum bad pair.
\end{enumerate}
\end{lemma}
\begin{proof}
The properness clause of \noderef{TashTreeMinimumBadPair} and
\noderef{MissingColorIncidenceEndpoint} show that the vertex-node \(y_j\)
has degree at most one in the two-color incidence graph, because it misses
\(\beta\).  Hence its component \noderef{TwoColorComponent} is a
\noderef{TwoColorPathComponent}; concretely,
\noderef{TwoColorComponentSpanningPath} gives a spanning path for it.
The same endpoint lemma says that \(y_i\), if it belongs to this component,
is also an end because it misses \(\alpha\).

Suppose \(y_i\) is not the other end.  It then does not belong to \(Q\):
if it did, it would be a second end of the spanning path, whose first end is
\(y_j\).  Apply the exact component swap
\noderef{TwoColorComponentSwap}.  It is proper by
\noderef{TwoColorComponentSwapProper}.  At \(y_j\), swapping the unique
possible incident \(\alpha\)-edge to \(\beta\) makes \(\alpha\) missing;
if the component is isolated, both colors were already missing and the
same conclusion holds.  Since \(y_i\notin Q\), \(\alpha\) remains missing
there.

Let \(T_j\) be the prefix through \(y_j\).  The unused hypotheses
\noderef{TreeColorUnusedThrough}, together with the fact that \(e_1=e\) is
uncolored, imply that no exact edge identity of \(T_j\) belongs to \(Q\).
Thus all its displayed edge colors stay fixed.  At a vertex of the prefix
strictly inside \(Q\), both colors are present before and after the swap;
at a vertex outside \(Q\), its missing colors do not change.  The only new
endpoint missingness in the prefix is the missing \(\alpha\) at \(y_j\).
It follows that every old Tashkinov witness in the prefix persists, while
\(\alpha\) is now missing at both \(y_i\) and \(y_j\).  The prefix is a
non-elementary Tashkinov tree, contradicting
\noderef{TashTreeProperPrefixElementary}.  Therefore \(Q\) joins
\(y_j\) to \(y_i\), proving the first assertion.

Now perform the swap on this path.  At its end \(y_i\), the possible
incident \(\beta\)-edge becomes \(\alpha\), so \(\beta\) is missing; at
the other end \(y_j\), the possible \(\alpha\)-edge becomes \(\beta\),
so \(\alpha\) is missing.  This also covers a degree-zero end directly.
The absence of both colors from \(e_2,\ldots,e_j\) proves the asserted
edge-by-edge fixity through the prefix.

We verify every later witness explicitly.  Fix a later displayed edge
\(e_h\) and an old witness \(y_r\), \(r<h\), for its old color.  A color
outside \(\{\alpha,\beta\}\) is untouched on both the edge and at the
witness.  If the old color is \(\alpha\) or \(\beta\) and \(y_r\) is not
one of \(y_i,y_j\), then a missing affected color makes \(y_r\) a path
end by \noderef{MissingColorIncidenceEndpoint}; since the two ends of
\(Q\) are already \(y_i,y_j\), either \(y_r\notin Q\) and its missingness
is unchanged, or it equals one of those two vertices.  If the witness
\(y_i\) for \(\alpha\) is lost, use \(y_j\), which is earlier than
\(e_h\) and now misses \(\alpha\).  If the witness \(y_j\) for \(\beta\)
is lost, use the earlier \(y_i\), which now misses \(\beta\).  Finally, if
\(e_h\) itself belongs to \(Q\), then an old \(\alpha\)-edge becomes
\(\beta\) and has the new earlier witness \(y_i\), while an old
\(\beta\)-edge becomes \(\alpha\) and has the new earlier witness
\(y_j\).  These are exactly the pointwise witness alternatives asserted
in the statement.  They prove that \(T\) remains a Tashkinov tree.

The exact uncolored edge \(e\) is not a node of the two-color incidence
graph, so the swap fixes it and fixes the property that it is the unique
uncolored edge.  All graph and numerical clauses of the \(k\)-triple are
unchanged, and properness was proved above.  At internal vertices of \(Q\)
both affected colors remain present; outside \(Q\) nothing changes; and at
the two ends the missing statuses of \(\alpha,\beta\) are exchanged.
Therefore the number of tree vertices missing each of the two colors is
unchanged, as is every other missing-color multiplicity.  Since \(T\) was
non-elementary, it remains non-elementary.  Thus
\((\varphi',T)\) is bad.  Its displayed sequence, path index, and rank are
unchanged.  Moreover the minimality clause of
\noderef{TashTreeMinimumBadPair} already excludes smaller bad pairs for
every generating coloring, so it applies unchanged to \(\varphi'\).
This proves the final assertion.
\end{proof}
